% Generated draft: requires mathematical and source review.
\documentclass[11pt]{article}
\usepackage[margin=25mm]{geometry}
\usepackage{fontspec}
\setmainfont{Latin Modern Roman}
\usepackage{amsmath,amssymb,mathtools,unicode-math}
\setmathfont{Latin Modern Math}
\usepackage{enumitem,needspace,xurl,hyperref}
\setlength{\parindent}{0pt}
\setlength{\parskip}{5pt}
\setlength{\emergencystretch}{3em}
\newcommand{\N}{\mathbb{N}}
\newcommand{\Z}{\mathbb{Z}}
\newcommand{\Q}{\mathbb{Q}}
\newcommand{\R}{\mathbb{R}}
\newcommand{\C}{\mathbb{C}}
\newcommand{\F}{\mathbb{F}}
\newcommand{\A}{\mathbb{A}}
\newcommand{\PP}{\mathbb{P}}
\newcommand{\E}{\mathbb{E}}
\newcommand{\eps}{\varepsilon}
\newcommand{\dd}{\,\mathrm{d}}
\begin{document}
\section*{UnsolvedMath ID: 6700050}
\textbf{Source problem code:} catalogue entry
\subsection{Definitions and mathematical statement}
Let \(S^n\) denote the unit round sphere, whose scalar curvature is \(n(n-1)\). Gromov's \emph{Sharp Spherical Length Comparison Inequality} is presented in three related forms:
\begin{enumerate}
\item Spheres with finitely many punctures are length extremal.
\item Every Riemannian \(n\)-manifold \(X\), possibly incomplete and possibly with boundary, with \(\operatorname{Scal}_X\geq n(n-1)\) satisfies
\[
 \operatorname{co\text{-}s.leng}(X)\leq2\pi.
\]
\item In the source's plain-language formulation, no strictly distance decreasing proper map \(f:X\to S^n\) of nonzero degree exists under the same scalar curvature bound.
\end{enumerate}
The source calls the equivalence only ``probably'' true. Neither length extremality nor \(\operatorname{co\text{-}s.leng}\) is defined in the excerpt, so their precise metric and topological conventions are unresolved.

With ordinary conventions, proper means that inverse images of compact sets are compact, and strict distance decrease can mean \(d_{S^n}(f(x),f(y))<d_X(x,y)\) for all distinct \(x,y\); a uniform Lipschitz bound strictly below one is a stronger condition in general. The source does not specify which condition is intended. A degree for a map of closed connected oriented manifolds is characterized by \(f_*[X]=(\deg f)[S^n]\), but the stated incomplete and boundary cases require a suitable degree convention not provided here.

There is also a target issue: ordinary properness for a map to the compact sphere forces \(X\) compact. Thus the literal third form does not cover noncompact punctured spheres without an additional convention, such as a specified noncompact target or relative condition; none is silently supplied. No spin, completeness, or boundary mean curvature hypothesis is added to the proposed general conjecture.
\subsection{Short English statement}
Gromov proposes sharp length comparison with the sphere, including finitely punctured spheres and incomplete manifolds with boundary under the spherical scalar curvature lower bound. The asserted equivalence and the conventions for proper maps, degree, and the co-spherical length invariant need clarification.
\subsection{Sources}
\begin{itemize}
\item Local UnsolvedMath registry, permanent ID 6700050, source code catalogue entry; Misha Gromov, 2017; update 2026-08-06.
\item Gromov, \emph{101 Questions, Problems and Conjectures around Scalar Curvature} (2017), printed tag [?53], occurrence 52, PDF page 57, ``Sharp Spherical Length Comparison Inequality'': \url{https://www.ihes.fr/~gromov/wp-content/uploads/2018/08/101-problemsOct1-2017.pdf}. Registry access date 2026-07-29; normalized PDF extraction; public author-hosted incomplete PDF; current status and reuse terms require review. The trailing 54 in the extraction appears to be a footnote marker whose text is not supplied.
\item The supplied partial-progress summary mentions Llarull-type spherical comparison and Gromov--Lawson index arguments for complete spin settings, without an exact citation or full hypotheses. Those reports do not establish the full incomplete, boundary, and non-spin conjecture above.
\end{itemize}
\end{document}
